% !PRINTED 603
\weberpage{603}

Ist $Q'$ gleich $A$, so ist $Q$ entweder auch gleich $A$ oder gleich $S$. Ist $Q'$ aber $=1$, so enthält $Q$ ausser der Einheit keine Permutation der ersten Art. Sind also $\varkappa$ und $\lambda$ zwei verschiedene und von der Einheit verschiedene Permutationen von $Q$, so müssen $\varkappa^2$ und $\varkappa\lambda$ als von der ersten Art $=1$ sein, d.\,h. $\lambda$ muss $=\varkappa$ sein. Es kann also $Q$ höchstens eine von der identischen verschiedene Permutation $\varkappa$ enthalten. Da aber $Q$ ein Normaltheiler von $S$ sein soll, so muss für jede Permutation $\pi$ aus der symmetrischen Gruppe $\pi^{-1}\varkappa\pi = \varkappa$ sein, d.\,h. $\varkappa$ darf sich nicht ändern, wenn in seinen Cyklen irgend eine Permutation ausgeführt wird. Dies ist aber nur dann möglich, wenn überhaupt nur zwei Ziffern 1, 2 vorhanden sind und $\varkappa = (1,\,2)$ ist. Dann aber ist $1, (1,\,2)$ die ganze Gruppe $S$.

\begin{center}
\S.\,178.

\gesperrt{Nicht metacyklische Gleichungen im K\"orper der rationalen Zahlen.}
\end{center}

Durch den Satz des vorigen Paragraphen ist der Nachweis geführt, dass eine Gleichung $n^{\text{ten}}$ Grades, wenn $n$ grösser als 4 ist, nicht mehr algebraisch gelöst werden kann, wenn die Co\"efficienten als unabhängige Veränderliche betrachtet werden. Von grösserem Interesse noch ist aber die Frage, ob es in dem Körper der rationalen Zahlen Gleichungen $n^{\text{ten}}$ Grades giebt, die nicht algebraisch lösbar sind. Die Frage lässt sich noch etwas allgemeiner stellen, nämlich so, ob es ganzzahlige Gleichungen giebt, deren Gruppe die symmetrische Gruppe ist, die also nach einer früher erklärten Ausdrucksweise keinen Affect haben.

Bildet man die \gesperrt{Galois}'sche Resolvente $G(t)=0$ vom Grade $\Pi(n)$ einer allgemeinen Gleichung $n^{\text{ten}}$ Grades $f(x)=0$ mit unbestimmten Co\"efficienten $a$, so ist $G(t)$ eine ganze Function der Veränderlichen $t$ und $a$, welche sich nicht in Factoren zerlegen lässt, die wieder rationale Functionen von $t$ und $a$ sind. Substituirt man für die Variablen $a$ Grössen irgend eines Körpers $\Omega$, und wird dann $G(t)$ in diesem Körper reducibel, so hat die Gleichung $f(x)=0$ im Rationalitätsbereich $\Omega$ einen Affect. Unsere Frage kommt also darauf hinaus, ob man in $G(t)$ für die Variablen $a$ solche rationale Zahlen setzen kann, dass $G(t)$ im Körper der rationalen Zahlen irreducibel bleibt.
